\documentclass[11pt,a4paper]{article}
\usepackage[utf8]{inputenc}
\usepackage[T1]{fontenc}
\usepackage[brazilian]{babel}
\usepackage{geometry}
\usepackage{helvet}
\usepackage[table]{xcolor}
\usepackage{enumitem}
\usepackage{array}
\usepackage{tabularx}
\renewcommand{\familydefault}{\sfdefault}
\geometry{top=1.8cm,bottom=1.8cm,left=2cm,right=2cm}
\setlength{\parindent}{0pt}
\setlength{\parskip}{0.5em}
\definecolor{azul}{HTML}{1F4E79}
\definecolor{azulclaro}{HTML}{EAF2F8}
\definecolor{cinza}{HTML}{F4F6F7}
\newcommand{\termo}[2]{\vspace{0.25em}\noindent\colorbox{azulclaro}{\parbox{\dimexpr\linewidth-2\fboxsep}{\textbf{#1} --- #2}}\par\vspace{0.25em}}
\newcommand{\secao}[1]{\vspace{0.8em}\noindent\textcolor{azul}{\Large\textbf{#1}}\par\vspace{0.2em}\hrule\vspace{0.35em}}
\pagestyle{plain}

\begin{document}

\begin{center}
{\LARGE\textbf{Apostila de Estudo}}\\[0.15cm]
{\Large\textcolor{azul}{Mudança de Paradigma na Ciência}}\\[0.2cm]
\textit{Questão 4 da prova 2026.1 --- baseado no material ``Evolução do Pensamento Científico e o Conceito de Método''.}
\end{center}

\secao{A ideia central}

Para Thomas Kuhn, a ciência não avança apenas acumulando descobertas. Durante um período, os cientistas trabalham a partir de um mesmo modelo aceito. Quando esse modelo passa a falhar diante de problemas persistentes, pode ocorrer uma mudança profunda: um novo modelo reorganiza o campo de estudos. Isso é uma \textbf{mudança de paradigma}.

\termo{Paradigma}{Modelo ou padrão aceito pela comunidade científica. Ele orienta quais perguntas são relevantes, quais métodos são usados e quais resultados são esperados.}

\secao{As quatro etapas: o processo completo}

\begin{center}
\fcolorbox{azul}{cinza}{\parbox{0.9\linewidth}{\centering\textbf{Ciência normal} $\longrightarrow$ \textbf{Anomalias} $\longrightarrow$ \textbf{Crise} $\longrightarrow$ \textbf{Revolução científica / novo paradigma}}}
\end{center}

\termo{1. Ciência normal}{É o trabalho científico realizado dentro do paradigma vigente. Os pesquisadores confiam no modelo aceito e procuram resolver problemas usando suas regras, teorias e técnicas.}

\termo{2. Anomalias}{São resultados, fatos ou problemas que não correspondem ao que o paradigma prevê. No início, podem parecer erro de medida, falha experimental ou detalhe ainda não explicado.}

\termo{3. Crise}{Acontece quando as anomalias persistem e não podem mais ser ajustadas ao modelo antigo. A confiança no paradigma diminui, surgem versões concorrentes e começam investigações extraordinárias.}

\termo{4. Revolução científica}{Ocorre quando um novo paradigma substitui total ou parcialmente o anterior. Mudam a forma de compreender o objeto, as técnicas e os compromissos da comunidade científica. Depois disso, inicia-se uma nova ciência normal.}

\secao{Não confunda os termos}

\begin{tabularx}{\linewidth}{>{\bfseries}p{3.1cm}X}
\hline
\rowcolor{azulclaro}Termo & O que ele significa \\
\hline
Paradigma & O modelo aceito que guia a ciência em um período. \\
Anomalia & Um problema que o modelo não explica adequadamente. \\
Crise & O enfraquecimento do modelo porque as anomalias se acumulam. \\
Revolução & A substituição do modelo antigo por outro. \\
Adaptação & Ajuste dentro do paradigma existente; não é necessariamente uma revolução. \\
\hline
\end{tabularx}

\secao{Exemplo do material}

O material cita a passagem da mecânica newtoniana para a mecânica relativística. Observações que entravam em conflito com expectativas do modelo anterior contribuíram para uma crise. A teoria da relatividade introduziu uma nova compreensão, substituindo a ideia de tempo absoluto pela de tempo relativo. O ponto do exemplo não é decorar datas: é reconhecer a sequência \textit{modelo aceito $\rightarrow$ dificuldade persistente $\rightarrow$ nova teoria $\rightarrow$ novo modelo}.

\secao{Como responder à prova}

Use esta estrutura de quatro movimentos:

\begin{enumerate}[leftmargin=0.7cm,itemsep=0.35em]
\item Defina paradigma como o modelo aceito pela comunidade científica.
\item Explique que a ciência normal trabalha conforme esse modelo.
\item Mostre que anomalias persistentes geram crise, pois o paradigma já não explica satisfatoriamente os resultados.
\item Conclua que a crise pode levar a uma revolução científica e à adoção de um novo paradigma.
\end{enumerate}

\textbf{Resposta-modelo:} Segundo Kuhn, a ciência inicialmente se desenvolve em uma fase de ciência normal, orientada por um paradigma aceito pela comunidade científica. Quando surgem anomalias, isto é, resultados que não correspondem às expectativas desse paradigma, tenta-se inicialmente explicá-las ou ajustá-las. Se elas persistem e se acumulam, instala-se uma crise. A crise estimula novas teorias e investigações, podendo culminar em uma revolução científica, na qual um novo paradigma substitui total ou parcialmente o anterior e passa a orientar uma nova fase de ciência normal.

\secao{Variações de pergunta: como adaptar a resposta}

\begin{tabularx}{\linewidth}{>{\bfseries}p{4.6cm}X}
\hline
\rowcolor{azulclaro}Se a pergunta pedir... & A ideia que não pode faltar \\
\hline
``O que é uma anomalia?'' & É um resultado que não se encaixa nas expectativas do paradigma vigente. \\
``Por que ocorre uma crise?'' & Porque anomalias persistentes não conseguem ser resolvidas pelo modelo antigo. \\
``Toda anomalia gera revolução?'' & Não. Algumas são resolvidas por adaptação; a revolução exige uma mudança de paradigma. \\
``Por que não é só acúmulo de conhecimento?'' & Porque um novo paradigma reorganiza a compreensão do campo, e não apenas acrescenta uma informação ao modelo anterior. \\
\hline
\end{tabularx}

\secao{Teste de domínio}

Você domina o tema se consegue responder, sem consultar:

\begin{itemize}[leftmargin=0.7cm,itemsep=0.25em]
\item Por que uma anomalia isolada ainda não é uma revolução científica?
\item Qual é a diferença entre anomalia e crise?
\item O que muda quando um novo paradigma é aceito?
\item Por que Kuhn diz que o avanço científico não é apenas acúmulo de conhecimentos?
\end{itemize}

\vfill
\small\textcolor{azul}{Atalho mental: modelo aceito $\rightarrow$ problema que não encaixa $\rightarrow$ perda de confiança $\rightarrow$ novo modelo.}

\end{document}
